\chapter{积分方程、Fredholm 理论与逆问题}
\label{chap:integral-equations}

\section{线性积分方程的分类}

积分方程的未知量是一个函数，而未知函数出现在积分号内。把最常见的一维线性情形写成

\[
\alpha(x)y(x)
=\phi(x)+\lambda\int_{a(x)}^{b(x)}K(x,t)y(t)\,\dd t.
\]

这里 $y$ 是未知函数，$\phi$ 是已知函数，$K(x,t)$ 称为积分核，$\lambda$ 是参数。这个式子先给出一般结构，具体分类只是对 $\alpha$ 和积分区间作特殊化。

若上下限都是固定的，称为 Fredholm 型；若至少有一个积分限随 $x$ 变化，最典型的是 $\int_a^x$，称为 Volterra 型。若 $\alpha(x)\equiv0$，未知函数只出现在积分号内，是第一类方程；若 $\alpha(x)$ 处处不为零，可以同除以 $\alpha(x)$ 化成第二类方程。

因此四种标准形式分别是

\[
\phi(x)=\int_a^bK(x,t)y(t)\,\dd t
\qquad\text{(Fredholm 第一类)},
\]

\[
y(x)=\phi(x)+\lambda\int_a^bK(x,t)y(t)\,\dd t
\qquad\text{(Fredholm 第二类)},
\]

\[
\phi(x)=\int_a^xK(x,t)y(t)\,\dd t
\qquad\text{(Volterra 第一类)},
\]

\[
y(x)=\phi(x)+\lambda\int_a^xK(x,t)y(t)\,\dd t
\qquad\text{(Volterra 第二类)}.
\]

若 $\alpha(x)$ 有孤立零点，就不能简单地处处同除以 $\alpha$，必须单独分析这些零点附近的可解性。

积分方程之所以在数学物理中自然出现，一个重要原因是微分算符的逆常常由 Green 函数给出。设边值问题

\[
Ly=f
\]

在给定边界条件下存在逆算符，并令 Green 函数满足

\[
L_xG(x,t)=\delta(x-t)
\]

以及相同的齐次边界条件，那么形式精确解为

\[
\boxed{
y(x)=\int G(x,t)f(t)\,\dd t
}.
\]

如果原问题是本征值方程

\[
Ly=\lambda w(x)y,
\]

并且 $0$ 不属于 $L$ 在这些边界条件下的谱，使 $L^{-1}$ 确实存在，就得到

\[
\boxed{
y(x)=\lambda\int G(x,t)w(t)y(t)\,\dd t
}.
\]

这说明 Sturm--Liouville 本征值问题可以转化成齐次 Fredholm 第二类方程。若 $L$ 有零模，Green 算符不再是普通意义下的全空间逆，此时必须先投影掉零模或加入相应正交可解条件。

具体地，Sturm--Liouville 算符

\[
L y=-\frac{\dd}{\dd x}
\left[p(x)\frac{\dd y}{\dd x}\right]+q(x)y
\]

在 $a<x<b$ 上配以例如

\[
y(a)\cos\alpha-y'(a)\sin\alpha=0,
\]

\[
y(b)\cos\beta+y'(b)\sin\beta=0.
\]

若相应 Green 函数存在，就有上面的积分表示。这一表示比先猜具体本征函数更一般，因为它直接把微分算符的边界条件编码进核 $G$ 中。

初值问题则自然产生 Volterra 方程。例如

\[
y'(x)=f(x,y),\qquad y(x_0)=y_0
\]

与积分方程

\[
\boxed{
y(x)=y_0+\int_{x_0}^x f(s,y(s))\,\dd s
}
\]

等价。二阶初值问题

\[
y''(x)=f(x,y),
\qquad y(x_0)=y_0,
\qquad y'(x_0)=y_1
\]

连续积分两次得到

\[
\boxed{
y(x)=y_0+y_1(x-x_0)
+\int_{x_0}^x(x-s)f(s,y(s))\,\dd s
}.
\]

这里核 $x-s$ 是由两次积分精确产生的，并没有做近似。

RC 电路也是同一结构。若 $i(t)$ 为电流，电容电荷满足

\[
q(t)=q(0)+\int_0^t i(s)\,\dd s,
\]

因此 Kirchhoff 电压定律可以直接写成关于 $i$ 的 Volterra 积分方程。

第一类方程中另一个经典例子是 Abel 方程

\[
\int_0^h\frac{u(z)}{\sqrt{h-z}}\,\dd z=f(h),
\qquad h>0.
\]

它的核在 $z=h$ 处具有可积奇性，是后面分数积分思想的原型。

热方程也可用 Green 函数表示。对

\[
u_t-\kappa u_{xx}=0,
\qquad 0<x<l,
\]

配以齐次 Dirichlet 边界条件和初值 $u(x,0)=f(x)$，若因果 Green 函数满足

\[
(\partial_t-\kappa\partial_x^2)G
=\delta(t-t')\delta(x-x'),
\qquad G=0\quad(t<t'),
\]

并满足同样的空间边界条件，则

\[
\boxed{
u(x,t)=\int_0^lG(x,t;x',0)f(x')\,\dd x'
}.
\]

因此积分方程和 Green 函数不是两套彼此独立的方法：Green 函数就是微分算符逆的积分核。下一节将固定在 Fredholm 型方程上，用 Hilbert 空间与紧算符语言讨论其一般可解性。
Fredholm 方程的 Neumann 级数一般需要 $|\lambda|\|K\|<1$ 才能直接保证收敛，但 Volterra 算符因为积分区域具有三角结构，情况更强。考虑
\[
y(x)=f(x)+\lambda\int_a^xK(x,t)y(t)\,\dd t,
\qquad a\le x\le b,
\]
并假设连续核满足
\[
|K(x,t)|\le M.
\]
记
\[
(Vy)(x)=\int_a^xK(x,t)y(t)\,\dd t.
\]
第一次作用有
\[
|(Vy)(x)|\le M(x-a)\|y\|_\infty.
\]
第二次作用为
\[
(V^2y)(x)
=\int_a^x\!K(x,t_1)
\int_a^{t_1}\!K(t_1,t_2)y(t_2)\,\dd t_2\dd t_1.
\]
因此
\[
|V^2y(x)|
\le M^2\|y\|_\infty
\int_a^x\int_a^{t_1}\dd t_2\dd t_1
=\frac{M^2(x-a)^2}{2!}\|y\|_\infty.
\]
归纳得到
\begin{equation*}\boxed{
|V^ny(x)|
\le \frac{M^n(x-a)^n}{n!}\|y\|_\infty
}.\end{equation*}
于是
\[
\sum_{n=0}^\infty |\lambda|^n\|V^n f\|_\infty
\le \|f\|_\infty
\sum_{n=0}^\infty
\frac{(|\lambda|M(b-a))^n}{n!}
=\|f\|_\infty \ee^{|\lambda|M(b-a)}.
\]
所以 Volterra 第二类方程的 Neumann 级数
\begin{equation*}\boxed{
y=\sum_{n=0}^\infty \lambda^nV^nf}\end{equation*}
在有限区间上对任意有限 $\lambda$ 都一致收敛。

这个结论与 Fredholm 情形不同的根源不是 $V$ 的范数特别小，而是迭代积分的有序区域
\[
a<t_n<\cdots<t_2<t_1<x
\]
在 $n$ 维立方体中的体积只有
\[
\frac{(x-a)^n}{n!}.
\]
正是这个 $1/n!$ 压倒了任意固定的 $|\lambda|^n$。

唯一性也可直接证明。若 $y_1,y_2$ 是两个解，令 $z=y_1-y_2$，则
\[
|z(x)|\le |\lambda|M\int_a^x|z(t)|\,\dd t.
\]
由 Gronwall 不等式立刻得到 $z\equiv0$。因此在上述连续核条件下，Volterra 第二类方程不仅存在解，而且解唯一。
定义迭代核
\[
K_1(x,t)=K(x,t),
\]
\[
K_{n+1}(x,t)
=\int_t^xK(x,s)K_n(s,t)\,\dd s.
\]
由 $|K|\le M$ 可归纳得到
\[
|K_n(x,t)|
\le M^n\frac{(x-t)^{n-1}}{(n-1)!},
\qquad a\le t\le x\le b.
\]
因此级数
\[
R_\lambda(x,t)
=\sum_{n=1}^{\infty}\lambda^{n-1}K_n(x,t)
\]
对任意有限 $\lambda$ 一致收敛，并满足
\[
R_\lambda(x,t)
=K(x,t)+\lambda\int_t^xK(x,s)R_\lambda(s,t)\,\dd s.
\]
于是解可以写成
\begin{equation*}\boxed{
y(x)=f(x)+\lambda\int_a^xR_\lambda(x,t)f(t)\,\dd t
}.\end{equation*}
这与 Fredholm 预解核形式相似，但 Volterra 核的三角结构使其预解对整个复 $\lambda$ 平面都是整函数，不会像 Fredholm 情形那样在离散特征值处出现极点。
考虑经典 Abel 方程
\begin{equation}
\boxed{
f(x)=\int_0^x\frac{u(t)}{\sqrt{x-t}}\,\dd t
}.
\label{eq:abel-first-kind}
\end{equation}
它是第一类 Volterra 方程。核在 $t=x$ 处发散，但奇性是可积的。关键不是直接对含奇核的式子求导，而是再作用一次同类积分算符。

对 $x>0$，定义
\[
I(x)=\int_0^x\frac{f(s)}{\sqrt{x-s}}\,\dd s.
\]
代入 \eqnrefs{eq:abel-first-kind}：
\[
I(x)=\int_0^x\frac{\dd s}{\sqrt{x-s}}
\int_0^s\frac{u(t)}{\sqrt{s-t}}\,\dd t.
\]
在适当可积条件下交换积分次序，积分区域 $0\le t\le s\le x$ 改写为 $0\le t\le x$、$t\le s\le x$：
\[
I(x)=\int_0^xu(t)
\left[
\int_t^x\frac{\dd s}{\sqrt{(x-s)(s-t)}}
\right]\dd t.
\]
内层积分通过
\[
s=t+(x-t)\sin^2\theta
\]
化为
\[
\int_0^{\pi/2}2\,\dd\theta=\pi.
\]
所以
\[
I(x)=\pi\int_0^xu(t)\,\dd t.
\]
对 $x$ 求导，得到 Abel 反演公式
\begin{equation*}\boxed{
u(x)=\frac1\pi\frac{\dd}{\dd x}
\int_0^x\frac{f(s)}{\sqrt{x-s}}\,\dd s
}.\end{equation*}

这段推导说明了为什么 Abel 核与 ``半阶积分'' 有关。若定义 Riemann--Liouville 分数积分
\[
(I^\alpha f)(x)
=\frac1{\Gamma(\alpha)}
\int_0^x(x-t)^{\alpha-1}f(t)\,\dd t,
\qquad \alpha>0,
\]
那么 Abel 算子就是
\[
\sqrt\pi\,I^{1/2}.
\]
而上面的双重积分恰好体现半群性质
\[
I^{1/2}I^{1/2}=I^1.
\]
因此 Abel 反演不是孤立技巧，而是分数积分代数的最简单实例。
考虑
\[
-u''(x)=f(x),
\qquad 0<x<L,
\]
配以
\[
u(0)=u(L)=0.
\]
我们要构造 $G(x,s)$ 使
\[
-\partial_x^2G(x,s)=\delta(x-s),
\qquad
G(0,s)=G(L,s)=0.
\]
在 $x\neq s$ 处，$G$ 关于 $x$ 是线性函数，因此写成
\[
G(x,s)=
\begin{cases}
A(s)x,&x<s,\\
B(s)(L-x),&x>s.
\end{cases}
\]
连续性要求
\[
A(s)s=B(s)(L-s).
\]
在 $s$ 附近积分微分方程：
\[
-\bigl[G_x(s^+,s)-G_x(s^-,s)\bigr]=1.
\]
因为
\[
G_x(s^-,s)=A(s),
\qquad
G_x(s^+,s)=-B(s),
\]
故
\[
A(s)+B(s)=1.
\]
联立两式得到
\[
A(s)=\frac{L-s}{L},
\qquad
B(s)=\frac{s}{L}.
\]
所以
\begin{equation*}\boxed{
G(x,s)=
\begin{cases}
\dfrac{x(L-s)}{L},&0\le x\le s,\\[5pt]
\dfrac{s(L-x)}{L},&s\le x\le L.
\end{cases}
}\end{equation*}
于是
\[
u(x)=\int_0^LG(x,s)f(s)\,\dd s.
\]
这个例子把 Green 函数的四个构造条件全部展示出来：分段满足齐次方程、满足边界条件、在源点连续、导数具有由 $\delta$ 决定的跳跃。


\section{Fredholm 型方程}
\label{sec:fredholm}

这一节把第二类 Fredholm 方程写成算符形式

\[
f=\phi+\lambda Kf,
\qquad
(I-\lambda K)f=\phi.
\]

这样做的好处是：存在性、唯一性、逐次逼近、谱展开和第一类方程的不适定性都能放在同一套 Hilbert 空间语言中讨论。
在 $L^2(a,b)$ 上定义

\[
(Kf)(x)=\int_a^b K(x,y)f(y)\,\dd y.
\]

若

\[
\|K\|_{\mathrm{HS}}^2
=\int_a^b\int_a^b|K(x,y)|^2\,\dd x\,\dd y<\infty,
\]

则 $K$ 是 Hilbert--Schmidt 算符，从而必为紧算符。由 Cauchy--Schwarz 不等式，

\[
\begin{aligned}
|(Kf)(x)|^2
&\le \left(\int_a^b|K(x,y)|^2\,\dd y\right)
\left(\int_a^b|f(y)|^2\,\dd y\right).
\end{aligned}
\]

再对 $x$ 积分，得到

\[
\|Kf\|_2\le \|K\|_{\mathrm{HS}}\|f\|_2,
\]

因此

\[
\boxed{\|K\|_{\mathrm{op}}\le\|K\|_{\mathrm{HS}}}.
\]

一般不能把这里的不等号写成等号。

Hilbert--Schmidt 条件为什么还能推出紧性，可以通过有限秩逼近看清楚。由于

\[
K(x,y)\in L^2((a,b)\times(a,b)),
\]

可取有限和

\[
K_N(x,y)=\sum_{j=1}^{N}a_j(x)b_j^*(y)
\]

使

\[
\|K-K_N\|_{L^2((a,b)^2)}\to0.
\]

相应算符

\[
(K_Nf)(x)
=\sum_{j=1}^{N}a_j(x)\langle b_j|f\rangle
\]

的值域包含在有限维空间 $\operatorname{span}\{a_1,\ldots,a_N\}$ 中，所以 $K_N$ 是有限秩算符，从而紧。又由刚才的估计

\[
\|K-K_N\|_{\mathrm{op}}
\le\|K-K_N\|_{\mathrm{HS}}\to0.
\]

紧算符在算符范数极限下仍为紧算符，因此 $K$ 紧。这里的逻辑是

\[
\text{平方可积核}
\Rightarrow\text{有限秩核可控逼近}
\Rightarrow\text{紧性}.
\]

伴算符的核为

\[
K^\dagger(x,y)=K^*(y,x).
\]

所以当

\[
K(x,y)=K^*(y,x)
\]

时，$K$ 是自伴紧算符。
若

\[
|\lambda|\,\|K\|_{\mathrm{op}}<1,
\]

则

\[
(I-\lambda K)^{-1}
=\sum_{n=0}^{\infty}(\lambda K)^n
\]

在算符范数下收敛，因此方程有唯一解

\begin{equation*}\boxed{f=\sum_{n=0}^{\infty}\lambda^nK^n\phi}.\end{equation*}

这不是形式猜测。令

\[
S_N=\sum_{n=0}^N(\lambda K)^n,
\]

则

\[
(I-\lambda K)S_N=I-(\lambda K)^{N+1}.
\]

由 $\|\lambda K\|<1$ 可知余项范数

\[
\|(\lambda K)^{N+1}\|
\le\|\lambda K\|^{N+1}\to0.
\]

所以这个展开的控制参数就是 $|\lambda|\|K\|$。当它接近或超过 $1$ 时，Neumann 级数不能作为一般保证，但原 Fredholm 方程仍可能有解；此时要转向谱理论或 Fredholm 二择一定理。

逐次逼近

\[
f_0=\phi,\qquad f_{n+1}=\phi+\lambda Kf_n
\]

正是 Neumann 部分和的递推形式。

若把算符幂重新写成积分核，还能得到经典的 Fredholm 预解式 核。定义迭代核

\[
K_1(x,y)=K(x,y),
\]

\[
K_{n+1}(x,y)
=\int_a^bK(x,z)K_n(z,y)\,\dd z.
\]

于是

\[
(K^n\phi)(x)=\int_a^bK_n(x,y)\phi(y)\,\dd y.
\]

在 Neumann 收敛条件下定义

\[
R_\lambda(x,y)
=\sum_{n=1}^{\infty}\lambda^{n-1}K_n(x,y),
\]

则精确解可写为

\[
\boxed{
f(x)=\phi(x)+\lambda\int_a^bR_\lambda(x,y)\phi(y)\,\dd y
}.
\]

把级数向前错一项还得到 预解恒等式

\[
R_\lambda=K+\lambda KR_\lambda
=K+\lambda R_\lambda K.
\]

因此 Neumann 级数、逐次逼近和 预解式 核 不是三套不同方法，而是同一个逆算符 $(I-\lambda K)^{-1}$ 的三种写法。
对紧算符 $K$，固定 $\lambda\neq0$。非齐次方程

\[
(I-\lambda K)f=\phi
\]

可解，当且仅当 $\phi$ 与伴随齐次方程

\[
(I-\lambda^*K^\dagger)g=0
\]

的所有解正交：

\begin{equation*}\boxed{\langle g|\phi\rangle=0}.\end{equation*}

特别地，若 $K$ 自伴且 $\lambda$ 为实数，则条件化为：$\phi$ 必须与

\[
(I-\lambda K)g=0
\]

的所有齐次解正交。

若齐次方程只有零解，则 $I-\lambda K$ 可逆，因而对每个 $\phi$ 都有唯一解。这就是``二择一''的核心：要么齐次问题出现非零解，要么非齐次问题对任意右端都有唯一解。

\begin{derivation}[Fredholm 正交条件与闭值域]
正交可解条件本身来自一个一般恒等式。对稠密定义闭算符 $A$，有
\[
\operatorname{Ran}(A)^\perp=\ker(A^\dagger).
\]
确实，$g\perp\operatorname{Ran}(A)$ 等价于
\[
\langle g|Af\rangle=0,
\qquad \forall f\in D(A),
\]
按伴算符定义这恰好说明 $g\in D(A^\dagger)$ 且 $A^\dagger g=0$。因此只要 $\operatorname{Ran}(A)$ 是闭的，就有
\[
\operatorname{Ran}(A)=\ker(A^\dagger)^\perp.
\]
现在取
\[
A=I-\lambda K,
\]
其中 $K$ 紧。关键是证明 $A$ 的值域闭。先在 $\ker A$ 的正交补上证明存在常数 $c>0$ 使
\[
\|Af\|\ge c\|f\|,
\qquad f\perp\ker A.
\]
若不存在，就能取 $\|f_n\|=1$、$f_n\perp\ker A$ 而 $Af_n\to0$。紧性使 $Kf_n$ 有收敛子列；由
\[
f_n=Af_n+\lambda Kf_n
\]
可知同一子列的 $f_n$ 也收敛到某个 $f$。极限满足 $Af=0$ 且 $f\perp\ker A$，同时 $\|f\|=1$，矛盾。因此上述下界成立。

若 $Af_n$ 在值域中收敛，把 $f_n$ 分解为 $\ker A$ 分量和正交分量；去掉不影响 $Af_n$ 的核分量后，下界说明剩余部分为 Cauchy 序列，所以极限仍在值域中。故 $\operatorname{Ran}(A)$ 闭。

于是非齐次方程 $Af=\phi$ 的可解性等价于
\[
\phi\perp\ker A^\dagger,
\]
即前面的 Fredholm 正交条件。还剩下二择一定理中最关键的一步：为什么 $\ker A=\{0\}$ 时，值域一定就是整个 Hilbert 空间？这里不把它简单归结为“Fredholm 指数为零”，而直接用紧算符的有限秩逼近算一遍。
\end{derivation}

要完成 Fredholm 二择一，还需把“齐次方程只有零解”提升为满射性。紧性允许用有限秩算符逼近 $\lambda K$，从而把满射性先化为有限维问题，再用一个小范数扰动恢复原算符的可逆性。

\begin{derivation}[由单射推出满射的有限秩逼近]

若 $\ker A=\{0\}$，上面已经证明存在 $c>0$ 使
\[
\|Af\|=\|(I-\lambda K)f\|\ge c\|f\|,
\qquad f\in\mathcal H.
\]
紧算符 $\lambda K$ 可以在算符范数下由有限秩算符 $F$ 逼近。取 $F$ 使
\[
\|\lambda K-F\|<\frac c4.
\]
于是
\[
\|(I-F)f\|
\ge \|Af\|-\|(\lambda K-F)f\|
\ge\frac{3c}{4}\|f\|.
\]
所以 $I-F$ 单射。由于 $F$ 只有有限维值域，$I-F$ 的满射性可以完全降到有限维问题。记
\[
M=\operatorname{Ran}F.
\]
给定任意 $y\in\mathcal H$，尝试寻找形如 $x=y+z$、$z\in M$ 的原像。方程 $ (I-F)x=y$ 等价于
\[
(I-F)z=Fy.
\]
这里左右两边都属于 $M$。限制算符
\[
(I-F)|_M:M\to M
\]
是单射，因为若 $z\in M$ 且 $(I-F)z=0$，由 $I-F$ 在整个空间上的单射性就有 $z=0$。有限维空间上单射等价于满射，因此对每个 $Fy\in M$ 都存在唯一 $z\in M$ 解上式。故每个 $y$ 都有原像 $x=y+z$，所以 $I-F$ 在整个 $\mathcal H$ 上双射，而且由上式
\[
\|(I-F)^{-1}\|\le\frac{4}{3c}.
\]
现在分解
\[
A
=(I-F)\left[I-(I-F)^{-1}(\lambda K-F)\right].
\]
第二个括号中的扰动满足
\[
\|(I-F)^{-1}(\lambda K-F)\|
<\frac{4}{3c}\frac c4
=\frac13<1.
\]
因此它可以用 Neumann 级数求逆。两个因子都可逆，所以 $A=I-\lambda K$ 可逆。于是
\[
\ker A=\{0\}
\quad\Longrightarrow\quad
\operatorname{Ran}A=\mathcal H.
\]
这就给出了 Fredholm 二择一定理“齐次解只有零解 $\Rightarrow$ 非齐次问题对任意右端唯一可解”的自包含证明，同时也说明紧性在其中究竟用在何处：它保证了 $\lambda K$ 可以被有限秩算符在算符范数下逼近。
\end{derivation}

上面的证明属于一个更一般的抽象框架。设 $T:X\to Y$ 是 Banach 空间之间的有界线性算子。若
\begin{equation*}\dim\ker T<\infty,
\qquad
\operatorname{Ran}T\text{ 闭},
\qquad
\dim(Y/\operatorname{Ran}T)<\infty,\end{equation*}
则称 $T$ 为 Fredholm 算子，并定义 指标
\begin{equation*}\boxed{
\operatorname{ind}T
:=\dim\ker T-\dim\operatorname{coker}T.
}\end{equation*}
在 Hilbert 空间中，闭值域定理给出
\[
\operatorname{coker}T
\cong(\operatorname{Ran}T)^\perp
=\ker T^\dagger,
\]
所以
\begin{equation*}\operatorname{ind}T
=\dim\ker T-\dim\ker T^\dagger.\end{equation*}
Fredholm 择一 的真正本质正是：$I-\lambda K$ 不仅值域闭，而且是 指标 为零的 Fredholm 算子。

为什么“恒等式 + 紧”天然进入 Fredholm 类，由 Atkinson 定理 精确回答。

\begin{theorem}[Atkinson theorem]\label{thm:math-fredholm-atkinson}
对 Banach 空间上的有界算子 $T$，下列条件等价：
\begin{enumerate}[label=(\roman*)]
\item $T$ 是 Fredholm 算子；
\item 存在有界算子 $S$ 使
\begin{equation}
ST-I\in\mathcal K,
\qquad
TS-I\in\mathcal K,
\label{eq:math-fredholm-parametrix-mod-compact}
\end{equation}
其中 $\mathcal K$ 表示紧算子理想。
\end{enumerate}
\end{theorem}

因此 Fredholm 性只取决于算子在 Calkin 代数
\[
\mathcal B(X)/\mathcal K(X)
\]
中的像是否可逆。换言之，“模掉 紧 error 后可逆”就是 Fredholm 的内在定义。对
\[
T=I-K,
\qquad K\text{ compact},
\]
直接取 $S=I$ 就有
\[
ST-I=TS-I=-K\in\mathcal K,
\]
故 $I-K$ 自动 Fredholm。

Atkinson 定理 还立即给出 紧 扰动 稳定性：若 $T$ Fredholm、$C$ 紧，则
\begin{equation*}T+C\text{ 仍 Fredholm},
\qquad
\boxed{\operatorname{ind}(T+C)=\operatorname{ind}T.}\end{equation*}
指标 是整数且在 Fredholm 算子集合上局部常数，因此沿连续 Fredholm 路径不会变化。特别地，考虑
\[
T_t=I-t\lambda K,
\qquad0\le t\le1.
\]
每个 $T_t$ 都是 恒等式 加 紧，故 Fredholm；而
\[
\operatorname{ind}T_0=\operatorname{ind}I=0.
\]
所以
\begin{equation*}\boxed{
\operatorname{ind}(I-\lambda K)=0.
}\end{equation*}
这条整数不变量正是 Fredholm 二择一的最短结构证明：若
\[
\ker(I-\lambda K)=\{0\},
\]
则 指标 为零迫使 cokernel 维数也为零；又因值域闭，所以值域就是整个空间。

\begin{derivation}[Atkinson 观点怎样统一二择一与稳定性]
设 $T$ Fredholm。把定义域分解为
\[
X=\ker T\oplus X_1,
\]
再把值域空间分解为
\[
Y=\operatorname{Ran}T\oplus Y_0,
\]
其中 $Y_0$ 有限维。限制
\[
T_1:=T|_{X_1}:X_1\to\operatorname{Ran}T
\]
是 Banach 空间同构。定义 $S:Y\to X$ 为
\[
S(y_r+y_0):=T_1^{-1}y_r.
\]
于是 $ST$ 在 $X_1$ 上为恒等、在 $\ker T$ 上为零，因此
\[
I-ST=P_{\ker T}
\]
是有限秩；同理
\[
I-TS=P_{Y_0}
\]
也是有限秩。于是得到\eqnrefs{eq:math-fredholm-parametrix-mod-compact}。

反过来，若存在 $S$ 满足\eqnrefs{eq:math-fredholm-parametrix-mod-compact}，则 $ST=I+K_1$、$TS=I+K_2$ 都是 恒等式 加 紧。前面的紧性论证保证它们的核有限维且值域闭；进一步可推出 $T$ 的核与 cokernel 都有限维。因此 $T$ Fredholm。

现在若加 紧 扰动 $C$，同一个 $S$ 满足
\[
S(T+C)-I=(ST-I)+SC\in\mathcal K,
\]
\[
(T+C)S-I=(TS-I)+CS\in\mathcal K,
\]
所以 $T+C$ 仍 Fredholm。指标 不变性则由路径 $T+tC$ 与 Fredholm 集合中的局部常值性得到：$t\mapsto\operatorname{ind}(T+tC)$ 是连续整值函数，故为常数。由此，“有限维 缺陷 对 紧 扰动 稳定”不再是积分核问题的特殊现象，而是 Fredholm category 的一般拓扑性质。

举例：Hilbert--Schmidt 算子是紧算子。设 $K:L^2(X)\to L^2(Y)$ 的核 $k(x,y)$ 满足
\[
\|K\|_{\mathrm{HS}}^2=\int_{X\times Y}|k(x,y)|^2\,\dd x\,\dd y<\infty,
\]
则 $K$ 是紧算子，且 $\|K\|\le\|K\|_{\mathrm{HS}}$。对 $T=I+K$，Atkinson 定理给出 $\operatorname{ind}T=0$，即 Fredholm 二择一：$Tx=y$ 可解当且仅当 $y\perp\ker K^*$。

Fredholm 指标在量子场论中描述 拓扑 扇区：Dirac 算子 $D$ 在规范场 $A$ 下的指标
\[
\operatorname{ind}D_A=\dim\ker D_A-\dim\ker D_A^\dagger
\]
由 Atiyah--Singer 定理给出，是规范场拓扑荷（瞬子数）的局部表达式。Atkinson 的 拟逆 $S$ 正是 Green 函数 $D_A^{-1}$ 在非零谱上的正则化，紧扰动 $C$ 对应规范场的连续形变，指标不变性保证物理 扇区 在连续形变下守恒。
\end{derivation}

紧算子 的谱结构也可以由这个观点直接读出。若 $K$ 紧，则对任意 $\mu\neq0$，
\[
K-\mu I=-\mu\left(I-\mu^{-1}K\right)
\]
是 Fredholm 指标 $0$。若它不可逆，Fredholm 择一 强迫非零 核，因此 $\mu$ 必为 本征值。再结合 紧性 可证明：每个非零 本征值 的代数重数有限，而且非零谱只能离散分布，唯一可能的聚点是 $0$。这就是 Riesz--Schauder 理论 的核心结论：
\begin{equation*}\boxed{
\sigma(K)\setminus\{0\}
=\{\text{有限重数的孤立本征值}\},
\quad
\text{唯一可能聚点为 }0.
}\end{equation*}
因此紧算子虽然作用在无限维空间上，但其非零谱仍保留有限维矩阵的离散本征值图像。

这个结构还能沿复参数解析延拓。设 $U\subset\CC$ 连通开，$K(z)$ 是取值于 紧算子 的 解析 family，并令
\[
A(z)=I-K(z).
\]
解析 Fredholm 定理 断言：只有两种可能。
\begin{enumerate}[label=(\roman*)]
\item 对每个 $z\in U$，$A(z)$ 都不可逆；
\item 存在至少一个 $z_0$ 使 $A(z_0)$ 可逆，此时不可逆点在 $U$ 中离散，并且
\begin{equation*}A(z)^{-1}\end{equation*}
是 算符-取值 亚纯 函数；每个极点的 主 部分 都是有限秩算子。
\end{enumerate}

\begin{derivation}[解析 Fredholm 定理 为什么只剩有限维奇性]
固定一个可能不可逆点 $z_*$。因为 $A(z_*)$ Fredholm 指标 $0$，可作有限维 缺陷 分解
\[
X=N\oplus X_1,
\qquad
X=R\oplus Y_0,
\]
其中
\[
N=\ker A(z_*),
\qquad
R=\operatorname{Ran}A(z_*),
\qquad
\dim N=\dim Y_0<\infty.
\]
在 $X_1\to R$ 这块，$A(z_*)$ 可逆；解析扰动后，在 $z_*$ 的小邻域仍可逆。于是通过 块 Gaussian elimination，可以把 $A(z)$ 的可逆性归约为一个有限维矩阵
\[
M(z):N\to Y_0
\]
的可逆性。$M(z)$ 的矩阵元关于 $z$ 解析，所以
\[
\det M(z)
\]
是普通复解析函数。它若不恒为零，其零点只能离散；有限维逆矩阵由 adjugate/行列式 表示
\[
M(z)^{-1}=\frac{\operatorname{adj}M(z)}{\det M(z)},
\]
因此只产生有限阶极点。把消元块重新乘回，就得到 $A(z)^{-1}$ 的 Laurent 主 部分 只作用在有限维 缺陷 空间 上，也就是有限秩。

若局部 行列式 在每个邻域都恒为零，解析延拓迫使 $A(z)$ 在整个连通域都不可逆，于是得到第一种情形。故无限维解析问题的全部非可逆奇性，最终仍被压缩成有限维解析 行列式。

举例：紧算子的预解式。取 $A(z)=I-zK$，$K$ 为 Hilbert--Schmidt 积分算子，核 $k(x,y)\in L^2([0,1]^2)$。设 $K$ 有非零本征值 $\mu_1,\mu_2,\ldots$，则
\[
(I-zK)^{-1}
\]
在 $z=1/\mu_j$ 处有极点，主 部分 为
\[
\sum_{k=1}^{m_j}\frac{P_k}{(z-1/\mu_j)^k},
\]
其中 $m_j$ 为 $\mu_j$ 的代数重数，$P_k$ 为有限秩投影。这正是量子散射理论中 共振 的数学结构。

例如 $K=(-\dd^2/\dd x^2+\omega^2)^{-1}$（Dirichlet，$L^2[0,\pi]$）本征值 $\mu_n=1/(n^2+\omega^2)$，$(I-zK)^{-1}$ 在 $z=n^2+\omega^2$ 处有单极点，主部为投影到 $\sin(nx)$ 的秩一算子，由此给出热核 $\ee^{-tH}=\sum\ee^{-t(n^2+\omega^2)}P_n$。反之 Volterra 算子 $(Kf)(x)=\int_0^x K(x,t)f(t)\dd t$ 谱仅为 $\{0\}$，$(I-zK)^{-1}$ 在全平面解析——这是定理第一种情形的实例。

物理上，紧扰动 $H=H_0+V$ 的预解式在复平面亚纯，极点即束缚态（实轴）与共振态（下半平面）；统计力学中 转移矩阵 的极点给出配分函数的相变零点（Yang--Lee），光学中共振腔模态的 $Q$ 因子由极点虚部决定。
\end{derivation}

对固定 紧算子 $K$，取
\[
A(z)=I-zK.
\]
因为 $A(0)=I$ 可逆，解析 Fredholm 定理 必落在第二种情形。因此
\begin{equation*}\boxed{
(I-zK)^{-1}
\text{ 关于 }z\text{ 亚纯},
}\end{equation*}
其极点恰位于
\[
z=\mu^{-1},
\qquad
\mu\in\sigma(K)\setminus\{0\}.
\]
这把 Neumann 级数的局部收敛区域升级成全复平面的亚纯 continuation：Neumann 级数只在 $|z|\|K\|<1$ 直接收敛，但 预解式 本身可以越过这一圆盘继续解析，直到遇到由非零本征值决定的离散极点。积分方程里的“共振参数”因此就是 算符-取值 亚纯 预解式 的极点。
若核可写成有限秩形式

\[
K(x,y)=\sum_{j=1}^N a_j(x)b_j^*(y),
\]

则

\[
(Kf)(x)=\sum_{j=1}^N a_j(x)c_j,
\qquad
c_j=\int_a^b b_j^*(y)f(y)\,\dd y.
\]

把

\[
f=\phi+\lambda Kf
\]

代入 $c_i$ 的定义，得到有限维线性方程组

\[
c_i=d_i+\lambda\sum_{j=1}^N M_{ij}c_j,
\]

其中

\[
d_i=\int_a^b b_i^*(x)\phi(x)\,\dd x,
\qquad
M_{ij}=\int_a^b b_i^*(x)a_j(x)\,\dd x.
\]

因此

\[
(I-\lambda M)\mathbf c=\mathbf d.
\]

这说明``退化核方法''的本质是把无限维积分方程精确降为有限维矩阵问题，而不是近似。只有当一般核再用有限秩核逼近时，才真正引入截断误差。

若

\[
\|K-K_N\|_{\mathrm{HS}}\to0,
\]

则由于 $\|K-K_N\|_{\mathrm{op}}\le\|K-K_N\|_{\mathrm{HS}}$，也得到算符范数控制。这样有限秩逼近才是受控近似。

一个有限秩共振例子可以把 Fredholm 二择一的抽象条件直接化成代数方程。

抽象的 Fredholm 择一 最容易在共振参数处显得难懂。下面用两个可以完全算完的例子，把``齐次解出现''、``非齐次右端必须正交''与``解不唯一''三件事放在同一张计算链上。

先看区间 $[0,1]$ 上的秩一核
\[
K(x,y)=1.
\]
第二类方程为
\begin{equation*}f(x)-\lambda\int_0^1 f(y)\,\dd y=\phi(x).\end{equation*}
记
\[
c=\int_0^1f(y)\,\dd y,
\qquad
d=\int_0^1\phi(y)\,\dd y.
\]
原方程立即给出
\[
f(x)=\phi(x)+\lambda c.
\]
再对 $x$ 从 $0$ 到 $1$ 积分，得到唯一的标量方程
\[
c=d+\lambda c,
\]
即
\[
(1-\lambda)c=d.
\]
因此若 $\lambda\neq1$，
\[
c=\frac{d}{1-\lambda},
\]
并且
\begin{equation*}\boxed{
f(x)=\phi(x)+\frac{\lambda}{1-\lambda}
\int_0^1\phi(y)\,\dd y
}.\end{equation*}
这里的预解核甚至可以直接读出：
\[
R_\lambda(x,y)=\frac1{1-\lambda}.
\]

真正有结构的是 $\lambda=1$。此时齐次方程
\[
f(x)-\int_0^1f(y)\,\dd y=0
\]
要求 $f$ 为任意常数函数，所以
\[
\ker(I-K)=\operatorname{span}\{1\}.
\]
由于该核是实对称核，$K^\dagger=K$，伴随齐次核也是同一空间。Fredholm 正交条件于是成为
\[
\langle1,\phi\rangle
=\int_0^1\phi(x)\,\dd x=0.
\]
这与标量方程在 $\lambda=1$ 时要求 $d=0$ 完全一致。若兼容条件满足，则
\[
f(x)=\phi(x)+c,
\qquad c\in\mathbb C
\]
全部都是解；若不满足，则根本无解。一个无限维积分方程在这个退化核上把 Fredholm 择一 的三种情况完整展示出来了。

第二个例子把积分方程和 Sturm--Liouville 共振直接连起来。考虑
\[
-u''(x)=g(x),
\qquad 0<x<\pi,
\qquad u(0)=u(\pi)=0.
\]
对应 Green 函数为
\[
G(x,y)=
\begin{cases}
\displaystyle \frac{x(\pi-y)}{\pi},&x\le y,\\[1ex]
\displaystyle \frac{y(\pi-x)}{\pi},&y\le x.
\end{cases}
\]
它满足
\[
-\partial_x^2G(x,y)=\delta(x-y),
\qquad G(0,y)=G(\pi,y)=0.
\]
因此逆算符
\[
(Kg)(x)=\int_0^\pi G(x,y)g(y)\,\dd y
\]
正是 Dirichlet 拉普拉斯算子 $L=-\dd^2/\dd x^2$ 的逆 $K=L^{-1}$。

现在研究带参数的边值问题
\begin{equation}
-u''(x)-\lambda u(x)=f(x),
\qquad u(0)=u(\pi)=0.
\label{eq:sturm-liouville-resonance}
\end{equation}
把它写成
\[
Lu=f+\lambda u
\]
并作用 $K=L^{-1}$，得到
\[
u=Kf+\lambda Ku,
\]
也就是
\[
(I-\lambda K)u=Kf.
\]
所以这里的 Fredholm 参数 $\lambda$ 恰好就是微分算符 $L$ 的谱参数。

Dirichlet 拉普拉斯算子 的正交归一本征函数是
\[
\varphi_n(x)=\sqrt{\frac2\pi}\sin nx,
\qquad n=1,2,\dots,
\]
满足
\[
L\varphi_n=n^2\varphi_n.
\]
因此
\[
K\varphi_n=\frac1{n^2}\varphi_n.
\]
Green 核也可展开为
\begin{equation*}G(x,y)
=\sum_{n=1}^\infty
\frac{2}{\pi n^2}
\sin(nx)\sin(ny),\end{equation*}
这里的每个 $1/n^2$ 都是紧算符 $K$ 的本征值。

若 $\lambda\neq n^2$ 对所有 $n$ 成立，把
\[
u=\sum u_n\varphi_n,
\qquad
f=\sum f_n\varphi_n
\]
代入微分方程，逐模态得到
\[
(n^2-\lambda)u_n=f_n,
\]
所以
\begin{equation*}\boxed{
u(x)
=\sum_{n=1}^\infty
\frac{f_n}{n^2-\lambda}\varphi_n(x)
}.\end{equation*}
当 $\lambda$ 接近某个 $m^2$ 时，第 $m$ 个模态的放大因子 $1/(m^2-\lambda)$ 变大；这就是预解算符在本征值附近出现极点的具体形式。

若恰好 $\lambda=m^2$，齐次问题存在
\[
u_h(x)=C\sin(mx).
\]
为了得到非齐次问题的兼容条件，不必先引用任何抽象定理。把\eqnrefs{eq:sturm-liouville-resonance} 两边乘以 $\sin(mx)$ 并积分：
\[
\int_0^\pi[-u''(x)-m^2u(x)]\sin(mx)\,\dd x
=\int_0^\pi f(x)\sin(mx)\,\dd x.
\]
对第一项分部积分两次，利用 $u(0)=u(\pi)=0$ 以及 $\sin(m\cdot0)=\sin(m\pi)=0$，左边恰好为零。因此必须有
\begin{equation}
\boxed{
\int_0^\pi f(x)\sin(mx)\,\dd x=0
}.
\label{eq:sturm-liouville-compatibility}
\end{equation}
这就是 $f\perp\ker(L-m^2I)$。因为 $L$ 自伴，伴随齐次空间仍由 $\sin(mx)$ 张成，所以它与 Fredholm 正交条件完全相同。

若\eqnrefs{eq:sturm-liouville-compatibility}成立，除第 $m$ 模态外其余系数仍由
\[
u_n=\frac{f_n}{n^2-m^2},
\qquad n\neq m
\]
确定，而 $u_m$ 任意。因此通解为
\[
u(x)
=\sum_{n\neq m}
\frac{f_n}{n^2-m^2}\varphi_n(x)
+C\varphi_m(x).
\]
至此可以把抽象 Fredholm 择一 翻译成一句非常具体的话：\textbf{当参数碰到齐次本征频率时，外力不能含有该共振模态；若没有该模态，特解存在，但还可以任意叠加齐次本征函数。}


自伴紧核允许把 Fredholm 方程进一步写成正交谱展开。

若 $K$ 是紧自伴算符，则存在一组正交归一的非零本征函数 $\{\varphi_n\}$，满足

\[
K\varphi_n=\mu_n\varphi_n,
\qquad \mu_n\in\mathbb R,
\]

非零本征值只能以 $0$ 为聚点。把

\[
f=\sum_n f_n\varphi_n+f_0,
\qquad
\phi=\sum_n\phi_n\varphi_n+\phi_0,
\]

其中 $f_0,\phi_0\in\ker K$，代入第二类方程，逐个模态得到

\[
(1-\lambda\mu_n)f_n=\phi_n.
\]

所以当 $1-\lambda\mu_n\neq0$ 时

\[
\boxed{f_n=\frac{\phi_n}{1-\lambda\mu_n}}.
\]

当 $\lambda=1/\mu_n$ 时分母为零，这正是 Fredholm 择一 中齐次方程出现非零解的谱论版本。

非自伴紧算子则应使用奇异值与 Schmidt 分解，而不是强行套用自伴本征展开。

若 $K$ 不自伴，不能指望本征函数本身给出正交基。此时考虑正算符

\[
K^\dagger K,\qquad KK^\dagger.
\]

它们具有相同的非零谱。若

\[
K^\dagger K\,\psi_n=s_n^2\psi_n,
\qquad s_n>0,
\]

定义

\[
\varphi_n=\frac1{s_n}K\psi_n,
\]

则

\[
K\psi_n=s_n\varphi_n,
\qquad
K^\dagger\varphi_n=s_n\psi_n.
\]

这就是紧算符的奇异系统。于是对 $f$ 的有效部分，

\[
Kf=\sum_n s_n\langle\psi_n|f\rangle\varphi_n.
\]

第一类方程

\[
Ky=g
\]

若写成

\[
g=\sum_n g_n\varphi_n,
\]

形式解为

\[
y=\sum_n\frac{g_n}{s_n}\psi_n.
\]

因为紧算符的奇异值 $s_n\to0$，小的 $s_n$ 会把数据误差放大成 $1/s_n$。这说明第一类 Fredholm 方程常常不适定，也说明为什么实际反演问题需要截断奇异值、Tikhonov 正则化等额外控制。

奇异系统还给出第一类方程可解性的精确 Picard 判据。若

\[
g=g_0+\sum_ng_n\varphi_n,
\qquad g_0\in\ker K^\dagger,
\]

那么 $Ky=g$ 存在平方可积解，当且仅当

\begin{equation*}\boxed{
g_0=0,
\qquad
\sum_n\frac{|g_n|^2}{s_n^2}<\infty
}.\end{equation*}

这时最小范数解为

\[
y^\dagger=\sum_n\frac{g_n}{s_n}\psi_n.
\]

若观测数据含噪声 $g^\delta=g+\eta$，小奇异值方向的噪声系数会被放大为 $\eta_n/s_n$。Tikhonov 正则化不再直接除以 $s_n$，而是解

\[
(K^\dagger K+\alpha I)y_\alpha=K^\dagger g^\delta,
\qquad \alpha>0.
\]

在奇异基上逐模态得到

\[
\boxed{
y_\alpha
=\sum_n\frac{s_n}{s_n^2+\alpha}\,g_n^\delta\,\psi_n
}.
\]

滤波因子

\[
\frac{s_n}{s_n^2+\alpha}
\]

在 $s_n\gg\sqrt\alpha$ 时近似 $1/s_n$，而在 $s_n\ll\sqrt\alpha$ 时不再爆炸。于是 $\alpha$ 明确充当``稳定性与逼近误差之间的控制参数''；令 $\alpha\to0$ 的同时必须与数据噪声水平一起选择，不能机械地把正则化参数直接取零。


\begin{derivation}[Tikhonov 正则化的正规方程与谱滤波]
给定噪声数据 $g^\delta$，考虑严格凸泛函
\[
J_\alpha(y)
=\|Ky-g^\delta\|^2
+\alpha\|y\|^2,
\qquad \alpha>0.
\]
令 $h$ 为任意扰动方向。对实参数 $\varepsilon$ 展开
\[
J_\alpha(y+\varepsilon h)
=J_\alpha(y)
+2\varepsilon\operatorname{Re}
\langle K^*(Ky-g^\delta)+\alpha y,h\rangle
+O(\varepsilon^2).
\]
极小点必须使一阶变分对所有 $h$ 为零，因此
\[
(K^*K+\alpha I)y_\alpha=K^*g^\delta.
\]
由于
\[
\langle y,(K^*K+\alpha I)y\rangle
=\|Ky\|^2+\alpha\|y\|^2
\ge\alpha\|y\|^2,
\]
算符 $K^*K+\alpha I$ 有界可逆，故极小点唯一。

在奇异系统 $K\psi_n=s_n\varphi_n$ 中写
\[
g^\delta=g_0^\delta+\sum_n g_n^\delta\varphi_n,
\qquad g_0^\delta\in\ker K^*,
\]
则正规方程逐模态化为
\[
(s_n^2+\alpha)y_n=s_ng_n^\delta.
\]
因此
\[
y_\alpha
=\sum_n\frac{s_n}{s_n^2+\alpha}g_n^\delta\psi_n.
\]
这说明正则化不是在形式逆上任意加一个“小量”，而是由一个 强制 变分问题产生的谱滤波。若数据误差满足 $\|g^\delta-g\|\le\delta$，要获得一致收敛还需令 $\alpha=\alpha(\delta)\to0$ 且噪声放大受控，例如典型条件 $\delta/\sqrt{\alpha}\to0$。
\end{derivation}

Riesz--Schauder 理论说明紧扰动为何在谱结构上保留有限维矩阵的许多性质。

Fredholm 二择一定理背后的结构可以再推进一步。设 $K$ 是无限维 Banach/Hilbert 空间上的紧算符，$\mu\neq0$。若
\[
Kf=\mu f,
\]
则 $\mu$ 是非零本征值。紧性强迫这类非零谱具有几乎有限维的结构：每个非零本征空间有限维，非零谱点不能在非零位置聚集，唯一可能的聚点是 $0$。

先证明非零本征空间有限维。若
\[
E_\mu=\ker(K-\mu I)
\]
无限维，则可在其中选取单位正交序列 $e_n$。因为
\[
Ke_n=\mu e_n,
\]
有
\[
\|Ke_n-Ke_m\|
=|\mu|\|e_n-e_m\|
=\sqrt2|\mu|
\qquad(n\neq m).
\]
所以 $Ke_n$ 没有 Cauchy 子列，与紧性矛盾。因此 $E_\mu$ 必有限维。

再证明非零本征值不能在非零位置聚集。反设存在彼此不同的本征值 $\mu_n$，并且对某个 $\delta>0$ 有
\[
|\mu_n|\ge\delta
\qquad\text{对所有 }n.
\]
把前 $n$ 个本征空间的和记为
\[
M_n=E_{\mu_1}+\cdots+E_{\mu_n}.
\]
不同本征值对应的本征向量线性无关，因此
\[
M_1\subsetneq M_2\subsetneq\cdots.
\]
每个 $M_n$ 都是有限维闭子空间。由 Riesz 引理，可以在 $M_n$ 中取单位向量 $x_n$，使
\[
\operatorname{dist}(x_n,M_{n-1})>\frac12.
\]
由于 $M_n$ 在 $K$ 下不变，而且在商空间 $M_n/M_{n-1}$ 上 $K$ 的作用就是乘以 $\mu_n$，所以
\[
Kx_n-\mu_nx_n\in M_{n-1}.
\]
对 $m<n$，有 $Kx_m\in M_{n-1}$，从而存在 $y_{mn}\in M_{n-1}$ 使
\[
Kx_n-Kx_m=\mu_nx_n+y_{mn}.
\]
于是
\[
\|Kx_n-Kx_m\|
\ge |\mu_n|\,\operatorname{dist}(x_n,M_{n-1})
>\frac{\delta}{2}.
\]
因此序列 $Kx_n$ 没有 Cauchy 子列，与 $K$ 的紧性矛盾。故对任意 $\delta>0$，满足 $|\mu|\ge\delta$ 的非零本征值只能有有限多个；非零谱唯一可能的聚点只能是 $0$。

这就是为什么紧算符在非零谱上像有限维矩阵，而真正的无限维效应集中在 $0$ 附近。第一类积分方程的不适定性也正发生在这里：奇异值不断趋向 $0$，逆算符因而无法有界。

把参数复化后，解析 Fredholm 理论进一步把不可逆点组织成预解算符的孤立极点。

考虑
\[
A(\lambda)=I-\lambda K.
\]
当 $|\lambda|\|K\|<1$ 时，Neumann 级数给出
\[
A(\lambda)^{-1}
=\sum_{n=0}^{\infty}\lambda^nK^n.
\]
但对紧算符，逆算符的存在区间远比这个圆盘更大。只要 $1/\lambda$ 不是 $K$ 的非零本征值，$I-\lambda K$ 就可逆。于是可能的奇点只出现在
\[
\lambda_n=\frac1{\mu_n}
\]
这些离散点上。

在自伴紧算符情形，这一点可以直接从谱展开看出。若
\[
Kf=\sum_n\mu_n\langle e_n|f\rangle e_n,
\]
则
\begin{equation*}\boxed{
(I-\lambda K)^{-1}f
=P_0f+
\sum_n\frac{1}{1-\lambda\mu_n}
\langle e_n|f\rangle e_n
},\end{equation*}
其中 $P_0$ 投影到 $\ker K$。每个因子
\[
\frac1{1-\lambda\mu_n}
\]
在 $\lambda=1/\mu_n$ 处有一阶极点。于是 Fredholm 预解式 的奇异性与齐次方程出现非零解完全对应。

这也说明 $|\lambda|\|K\|<1$ 只是 Neumann 级数的一条方便充分条件，不是方程可解的必要条件。真正的障碍是离散的谱点，而不是整个单位圆外区域。

第一类方程的不适定性可由奇异值方向上的误差放大精确量化。

对紧算符奇异系统
\[
K\psi_n=s_n\phi_n,
\qquad s_n\downarrow0,
\]
精确数据
\[
g=Ky^\dagger
\]
具有系数
\[
g_n=s_n y_n^\dagger.
\]
若观测数据变为
\[
g^\delta=g+\eta,
\qquad \|\eta\|\le\delta,
\]
直接形式反演会给出
\[
y_{\mathrm{formal}}^\delta
=\sum_n\left(y_n^\dagger+\frac{\eta_n}{s_n}\right)\psi_n.
\]
即使 $\eta$ 很小，只要噪声在足够小的 $s_n$ 方向上有分量，$\eta_n/s_n$ 仍可任意大。这正是 Hadamard 意义下的不稳定性。

Tikhonov 解为
\[
y_\alpha^\delta
=(K^\dagger K+\alpha I)^{-1}K^\dagger g^\delta
=\sum_n\frac{s_n}{s_n^2+\alpha}g_n^\delta\psi_n.
\]
把 $g_n^\delta=s_ny_n^\dagger+\eta_n$ 代入，得到精确误差分解
\begin{equation*}\boxed{
y_\alpha^\delta-y^\dagger
=-\sum_n\frac{\alpha}{s_n^2+\alpha}y_n^\dagger\psi_n
+\sum_n\frac{s_n}{s_n^2+\alpha}\eta_n\psi_n
}.\end{equation*}
第一项是正则化偏差，第二项是数据噪声传播。

噪声滤波因子满足
\[
0\le \frac{s}{s^2+\alpha}\le\frac1{2\sqrt\alpha},
\]
因为函数 $s/(s^2+\alpha)$ 在 $s=\sqrt\alpha$ 处取最大值。因此
\[
\left\|
(K^\dagger K+\alpha I)^{-1}K^\dagger\eta
\right\|
\le\frac{\delta}{2\sqrt\alpha}.
\]
于是 $\alpha$ 不能比噪声水平下降得过快，否则稳定性会再次丧失。

另一方面，若精确解具有额外光滑性，例如满足源条件
\[
y^\dagger=(K^\dagger K)^\nu w,
\qquad 0<\nu\le1,
\]
则
\[
y_n^\dagger=s_n^{2\nu}w_n.
\]
偏差系数变为
\[
\frac{\alpha s_n^{2\nu}}{s_n^2+\alpha}w_n.
\]
令 $r=s_n^2/\alpha$，可写成
\[
\alpha^\nu\frac{r^\nu}{1+r}w_n,
\]
而 $r^\nu/(1+r)$ 在 $0<\nu\le1$ 上有界，所以
\[
\|y_\alpha-y^\dagger\|
\le C_\nu\alpha^\nu\|w\|.
\]
总误差因此受控于
\[
\boxed{
\|y_\alpha^\delta-y^\dagger\|
\lesssim \alpha^\nu+\frac{\delta}{\sqrt\alpha}
}.
\]
这把“$\alpha$ 是稳定性与逼近误差之间的权衡参数”变成了定量结论，而不是一句定性描述。

截断奇异值与 Tikhonov 正则化只是对同一奇异谱采用不同滤波函数。

另一种直接稳定化方法是只保留
\[
s_n\ge\tau
\]
的奇异方向，定义
\[
y_\tau^\delta
=\sum_{s_n\ge\tau}\frac{g_n^\delta}{s_n}\psi_n.
\]
这称为截断奇异值分解。它的滤波函数是
\[
q_\tau(s)=
\begin{cases}
1/s,&s\ge\tau,\\
0,&s<\tau.
\end{cases}
\]
Tikhonov 的滤波函数则是
\[
q_\alpha(s)=\frac{s}{s^2+\alpha}.
\]
前者是硬截止，后者是平滑衰减。两者都遵循同一原则：不能在 $s\to0$ 时继续使用无控制的 $1/s$。因此不同正则化方法的核心并不是“换一种求逆公式”，而是重新设计小奇异值方向上的频谱滤波。三种滤波函数的差别见\figref{fig:tikhonov-filter}。


\begin{figure}[htbp]
\centering
\includegraphics[width=0.76\textwidth]{tikhonov_filter.pdf}
\caption{形式逆、Tikhonov 与截断奇异值的滤波函数比较。形式逆 $1/s$ 在 $s\to0$ 时无界；Tikhonov 用平滑方式压制小奇异值，TSVD 则在阈值以下直接截断。}
\label{fig:tikhonov-filter}
\end{figure}



\section{不适定反演的统一谱滤波框架}
\label{sec:inverse-spectral-filter}

前面的奇异系统已经暴露了第一类方程的本质困难：在 $s_n\to0$ 的方向上，形式逆包含 $1/s_n$。Tikhonov 只是处理这一困难的一种方法。为了比较不同算法，应先把它们写成同一个 泛函-演算 形式。

设 $K:X\to Y$ 为 Hilbert 空间之间的紧算子，$A:=K^*K\ge0$。对一族标量函数 $g_\alpha:[0,\|K\|^2]\to\RR$，定义正则化算子
\[
\boxed{
R_\alpha
:=g_\alpha(A)K^*
},
\qquad
x_\alpha^\delta=R_\alpha y^\delta.
\]
若 $K\psi_n=s_n\varphi_n$，则
\[
x_\alpha^\delta
=\sum_n s_ng_\alpha(s_n^2)
\langle y^\delta,\varphi_n\rangle\psi_n.
\]
形式逆对应 $g(\lambda)=1/\lambda$，而正则化的任务正是让 $g_\alpha(\lambda)$ 在 $\lambda\approx0$ 附近保持可控，同时在固定 $\lambda>0$ 上随 $\alpha\to0$ 恢复 $1/\lambda$。

定义 残差 函数
\[
r_\alpha(\lambda)
:=1-\lambda g_\alpha(\lambda).
\]
对最小范数精确解 $x^\dagger$，误差可以精确分成
\[
\boxed{
x_\alpha^\delta-x^\dagger
=R_\alpha(y^\delta-y)
-r_\alpha(A)x^\dagger
}.
\]
第一项是噪声传播，第二项是正则化偏差。于是所有参数选择问题都归结为同一个竞争：$\alpha$ 太小会放大噪声，太大则不能充分恢复小奇异值方向。

若
\[
\|y^\delta-y\|\le\delta,
\]
则噪声项满足
\[
\|R_\alpha(y^\delta-y)\|
\le
\delta\,\|R_\alpha\|,
\]
而谱定理给出
\[
\boxed{
\|R_\alpha\|
=
\sup_{0\le\lambda\le\|K\|^2}
\sqrt\lambda\,|g_\alpha(\lambda)|
}.
\]
因此算法稳定性不是由名称决定，而是由这一 sup 范数 精确决定。
Tikhonov 正则化对应
\[
g_\alpha^{\mathrm{Tik}}(\lambda)
=\frac{1}{\lambda+\alpha},
\qquad
r_\alpha^{\mathrm{Tik}}(\lambda)
=\frac{\alpha}{\lambda+\alpha}.
\]
于是
\[
\|R_\alpha^{\mathrm{Tik}}\|
=\sup_{\lambda\ge0}\frac{\sqrt\lambda}{\lambda+\alpha}
=\frac{1}{2\sqrt\alpha}.
\]
这直接解释了为什么典型噪声项按 $\delta/\sqrt\alpha$ 放大。

截断奇异值（谱 截断）则取
\[
g_\alpha^{\mathrm{TSVD}}(\lambda)
=\begin{cases}
\lambda^{-1},&\lambda\ge\alpha,\\
0,&0\le\lambda<\alpha,
\end{cases}
\]
所以
\[
r_\alpha^{\mathrm{TSVD}}(\lambda)
=\mathbf 1_{[0,\alpha)}(\lambda),
\qquad
\|R_\alpha^{\mathrm{TSVD}}\|=\alpha^{-1/2}.
\]
它不是``近似地除以奇异值''，而是完全保留稳定模态、完全删除不稳定模态。

Landweber 迭代从
\[
x_{m+1}=x_m+\omega K^*(y^\delta-Kx_m),
\qquad
0<\omega<\frac{2}{\|K\|^2}
\]
出发。若 $x_0=0$，递推求和得到
\[
x_m
=\omega\sum_{j=0}^{m-1}(I-\omega A)^jK^*y^\delta,
\]
故它同样是谱滤波：
\[
\boxed{
g_m^{\mathrm{LW}}(\lambda)
=\frac{1-(1-\omega\lambda)^m}{\lambda}}.
\]
因此``迭代次数''本身就是正则化参数。若无限迭代到噪声被完全拟合，稳定性反而会丢失；这就是 iterative 正规化 中 early stopping 的数学来源。
仅知道 $x^\dagger\in X$ 时，通常只能证明 $x_\alpha^\delta\to x^\dagger$，却不能得到统一收敛速率。要得到 速率，需要描述真实解相对于 $K^*K$ 的光滑性。一个标准 源 条件 是
\[
\boxed{
x^\dagger=A^\nu w,
\qquad \nu>0,
\quad \|w\|\le C_w}.
\]
在奇异基上，这意味着
\[
\langle x^\dagger,\psi_n\rangle
=s_n^{2\nu}\langle w,\psi_n\rangle,
\]
即真实解在最不稳定的小奇异值方向上本来就具有额外衰减。

对 Tikhonov，偏差为
\[
x_\alpha-x^\dagger
=-r_\alpha(A)x^\dagger
=-\alpha(A+\alpha I)^{-1}A^\nu w.
\]
因此
\[
\|x_\alpha-x^\dagger\|
\le
\sup_{\lambda\ge0}
\frac{\alpha\lambda^\nu}{\lambda+\alpha}\,\|w\|.
\]
令 $\lambda=\alpha t$，得到
\[
\frac{\alpha\lambda^\nu}{\lambda+\alpha}
=
\alpha^\nu\frac{t^\nu}{1+t}.
\]
当 $0<\nu\le1$ 时，右侧 sup 有界，因此
\[
\boxed{
\|x_\alpha-x^\dagger\|
\le C_\nu\alpha^\nu\|w\|
}.
\]
与噪声项合并，得到
\[
\|x_\alpha^\delta-x^\dagger\|
\le
\frac{\delta}{2\sqrt\alpha}
+C_\nu\alpha^\nu\|w\|.
\]
平衡两项给出
\[
\alpha\asymp\delta^{\frac{2}{2\nu+1}},
\qquad
\boxed{
\|x_\alpha^\delta-x^\dagger\|
=O\!\left(
\delta^{\frac{2\nu}{2\nu+1}}
\right)}.
\]
这不是经验拟合，而是 稳定性 与 approximation 两个误差项的最优阶平衡。

普通 Tikhonov 的上述简单估计在 $\nu>1$ 时不再继续提高阶数，这体现了 资格/saturation：一个正则化方法能利用多少阶的 源 光滑性 是有限的。若真实解更光滑而算法已经饱和，需要 iterated Tikhonov、谱截断或其他更高 资格 的滤波器，不能指望把同一个 $\alpha$ 公式无限外推。

\begin{derivation}[一般滤波器的误差分解]
由 $y=Kx^\dagger$ 且 $A=K^*K$，有
\[
\begin{aligned}
x_\alpha^\delta-x^\dagger
&=g_\alpha(A)K^*y^\delta-x^\dagger\\
&=g_\alpha(A)K^*(y^\delta-y)
+g_\alpha(A)K^*Kx^\dagger-x^\dagger\\
&=g_\alpha(A)K^*(y^\delta-y)
+\bigl(A g_\alpha(A)-I\bigr)x^\dagger\\
&=R_\alpha(y^\delta-y)
-r_\alpha(A)x^\dagger,
\end{aligned}
\]
其中 $R_\alpha:=g_\alpha(A)K^*$ 为残差算子，$r_\alpha(\lambda):=1-\lambda g_\alpha(\lambda)$ 为残差滤波函数。由于 $K=U A^{1/2}$ 是 极 分解（$U$ 为部分等距，$\ker U=\ker K$），
\[
R_\alpha
=g_\alpha(A)A^{1/2}U^*,
\]
所以
\[
\|R_\alpha\|
\le
\|g_\alpha(A)A^{1/2}\|
=
\sup_{\lambda\in\sigma(A)}
\sqrt\lambda\,|g_\alpha(\lambda)|.
\]
在 $\overline{\operatorname{Ran}K}$ 上 $U^*$ 为部分等距，故该上界实际取等，由此即得精确公式。

两项的物理意义如下：第一项 $R_\alpha(y^\delta-y)$ 是噪声放大项，范数 $\|R_\alpha\|\delta$ 随 $\alpha\to0$ 发散（因 $g_\alpha(\lambda)\to1/\lambda$）；第二项 $r_\alpha(A)x^\dagger$ 是正则化偏差，随 $\alpha\to0$ 趋零。最优 $\alpha$ 平衡二者，给出相应的收敛率。

取 $g_\alpha(\lambda)=1/(\lambda+\alpha)$（Tikhonov），则 $\|R_\alpha\|=1/(2\sqrt\alpha)$、$\|r_\alpha(A)\|=1$；对 源 条件 $x^\dagger=A^\nu w$，平衡 $\delta/(2\sqrt\alpha)=\alpha^\nu\|w\|$ 得 $\alpha\asymp\delta^{2/(2\nu+1)}$。谱截断 $g_\alpha=\lambda^{-1}\mathbf 1_{\lambda\ge\alpha}$ 对紧算子给出 $x_\alpha^\delta=\sum_{s_j\ge\alpha}\langle y^\delta,v_j\rangle u_j/s_j$，$\|R_\alpha\|=1/\alpha$，同阶但对高频硬归零；Landweber 迭代 $g_k=\lambda^{-1}(1-(1-\lambda)^k)$ 给出 $k\asymp\delta^{-2/(2\nu+1)}$ 而无需显式谱分解；$m$ 次 iterated Tikhonov 把 资格 提到 $m$，突破普通 Tikhonov 的 saturation。

此谱滤波框架在各领域同构出现：量子逆问题中 $g_\alpha$ 是有限分辨率测量的分辨率函数，统计力学中与热浴 $\ee^{-\beta H}$（$\alpha\sim1/\beta$）对应，图像重建中 Tikhonov 对应 Tikhonov--Miller，CT/MRI 的滤波反投影即 Radon 变换的离散谱滤波。
\end{derivation}
若噪声水平 $\delta$ 已知，参数不能只根据图像``看起来最平滑''来选。Morozov discrepancy 原理 要求选取 $\alpha=\alpha(\delta,y^\delta)$ 使
\[
\boxed{
\|Kx_\alpha^\delta-y^\delta\|
\approx\tau\delta,
\qquad \tau>1
}.
\]
其逻辑是：残差若远小于噪声水平，说明算法已经开始拟合无法辨认的数据噪声；残差若远大于噪声水平，则说明正则化过强，仍有真实可辨认信息尚未被利用。

在谱滤波语言中，残差写成
\[
Kx_\alpha^\delta-y^\delta
=-r_\alpha(KK^*)y^\delta
\]
（在 $\ker K^*$ 方向还保留不可拟合分量）。因此 discrepancy 原理 实际上是在控制 残差 函数 对数据谱的作用。它把参数选择从``算法外部的经验旋钮''变成算子理论内部的一条停止准则。

\begin{note}[正则化不等于人为篡改方程]
原始方程 $Kx=y$ 在噪声数据上可能不存在稳定逆。正则化的目标不是宣称 $K$ 被另一个算子取代，而是构造一族连续近似逆 $R_\alpha$：固定 $\alpha>0$ 时稳定，随着数据误差 $\delta\to0$ 再让 $\alpha(\delta)\to0$，最终恢复最小范数解。两个极限必须联动；先在固定噪声下令 $\alpha\to0$ 通常恰好会重新暴露不适定性。
\end{note}

\section{Fourier 方法、卷积核与奇异积分}
\label{sec:integral-fourier}

当积分核具有平移不变性时，Fourier 变换会把积分算符对角化。这是上一节``谱展开''的连续谱版本。

先从一般方程

\[
y(s)=f(s)+\lambda\int_{-\infty}^{\infty}K(s,t)y(t)\,\dd t
\]

出发。采用\chapref{chap:distributions-fourier}的对称 Fourier 归一化，记

\[
Y(k)=\mathcal F[y](k),\qquad F(k)=\mathcal F[f](k).
\]

把 $y(t)$ 的逆 Fourier 表示代入，得到

\[
Y(k)=F(k)+\lambda\int_{-\infty}^{\infty}
\overline K(k,k')Y(k')\,\dd k',
\]

其中

\[
\boxed{
\overline K(k,k')
=\frac1{2\pi}
\int_{\mathbb R^2}
 \ee^{-\ii ks}K(s,t)\ee^{\ii k't}\,\dd s\,\dd t
}.
\]

这只是把同一个算符改写到 Fourier 表象中。只有在 $\overline K$ 近似或恰好对角时，方程才真正简化。

最重要的精确情形是平移不变核

\[
K(s,t)=\frac1{\sqrt{2\pi}}v(s-t).
\]

按照\chapref{chap:distributions-fourier}采用的归一化卷积

\[
(v*y)(s)=\frac1{\sqrt{2\pi}}
\int_{-\infty}^{\infty}v(s-t)y(t)\,\dd t,
\]

方程就是

\[
y=f+\lambda(v*y).
\]

由于

\[
\mathcal F[v*y]=VY,
\qquad V=\mathcal F[v],
\]

所以

\[
(1-\lambda V(k))Y(k)=F(k).
\]

只要分母不为零，就得到形式精确解

\[
\boxed{
Y(k)=\frac{F(k)}{1-\lambda V(k)}
},
\]

再作逆 Fourier 变换即可恢复 $y$。当 $1-\lambda V(k)$ 在某些 $k$ 上为零时，必须重新讨论谱点、边界值或广义函数意义下的可解性，不能直接把这个商当作普通函数。
奇异积分中最典型的是 Cauchy 主值。定义 Hilbert 变换

\[
\boxed{
(\mathcal H f)(x)
=\frac1\pi\operatorname{pv}
\int_{-\infty}^{\infty}\frac{f(t)}{x-t}\,\dd t
}.
\]

其核不是普通可积函数，而是主值分布

\[
\frac1\pi\operatorname{pv}\frac1x.
\]

若 $f(z)$ 在上半平面解析并具有足够快的无穷远衰减，对实轴点 $x$ 绕开一个上半平面小半圆。Cauchy 积分定理给出

\[
0=
\operatorname{pv}\int_{-\infty}^{\infty}
\frac{f(t)}{t-x}\,\dd t
-\ii\pi f(x),
\]

因此

\[
\boxed{
f(x)=\frac1{\ii\pi}
\operatorname{pv}\int_{-\infty}^{\infty}
\frac{f(t)}{t-x}\,\dd t
}.
\]

写

\[
f=u+iv,
\]

比较实部和虚部，得到一对共轭关系

\[
u(x)=\frac1\pi\operatorname{pv}
\int_{-\infty}^{\infty}\frac{v(t)}{t-x}\,\dd t,
\]

\[
v(x)=\frac1\pi\operatorname{pv}
\int_{-\infty}^{\infty}\frac{u(t)}{x-t}\,\dd t
=(\mathcal Hu)(x).
\]

在 Fourier 空间中，Hilbert 变换的乘子为

\[
\boxed{
\mathcal F[\mathcal H f](k)
=-\ii\,\operatorname{sgn}(k)\widehat f(k)
}.
\]

于是对没有 $k=0$ 奇异分量、并在相应函数空间中定义良好的 $f$，

\[
\mathcal H^2f=-f.
\]

这个公式说明主值积分算符同样可以通过 Fourier 谱乘子来理解。

上面从主值积分和解析函数边界值得到了 Hilbert 变换。若进一步把它放回 Hilbert 空间的语言中，几个常用性质可以一次性由 Fourier 乘子推出，而且这一步也把\chapref{chap:hilbert-space}的谱观点与本章的奇异积分真正接起来。

先在 Schwartz 空间 $\mathcal S(\mathbb R)$ 上定义
\[
(\mathcal H f)(x)=\frac1\pi\operatorname{pv}
\int_{-\infty}^{\infty}\frac{f(t)}{x-t}\,\dd t.
\]
按照\chapref{chap:distributions-fourier}的对称 Fourier 约定，已经有
\begin{equation*}\widehat{\mathcal H f}(k)
=-\ii\,\operatorname{sgn}(k)\widehat f(k).\end{equation*}
右端的乘子 $-\ii\operatorname{sgn}(k)$ 在几乎处处意义下模长为 $1$。因此 Plancherel 定理直接给出
\[
\|\mathcal H f\|_2^2
=\int_{\mathbb R}
\bigl|-\ii\operatorname{sgn}(k)\widehat f(k)\bigr|^2\,\dd k
=\int_{\mathbb R}|\widehat f(k)|^2\,\dd k
=\|f\|_2^2.
\]
所以 $\mathcal H$ 在 $\mathcal S(\mathbb R)$ 上是等距的。由于 $\mathcal S(\mathbb R)$ 在 $L^2(\mathbb R)$ 中稠密，$\mathcal H$ 唯一连续延拓为 $L^2(\mathbb R)$ 上的有界算符，而且
\begin{equation*}\|\mathcal H\|_{L^2\to L^2}=1.\end{equation*}
这说明主值积分虽然在 $x$ 空间中带有不可积奇点，但从 $L^2$ 算符的角度看却完全受控。

再看伴随。对 $f,g\in\mathcal S(\mathbb R)$，利用 Plancherel 定理和本书内积第一变量共轭线性的约定，
\begin{align*}
\langle \mathcal H f,g\rangle
&=\int_{\mathbb R}
\overline{-\ii\operatorname{sgn}(k)\widehat f(k)}\,
\widehat g(k)\,\dd k\\
&=\int_{\mathbb R}
\overline{\widehat f(k)}\,
\ii\operatorname{sgn}(k)\widehat g(k)\,\dd k\\
&=\langle f,-\mathcal H g\rangle.
\end{align*}
于是
\begin{equation}
\mathcal H^\dagger=-\mathcal H.
\label{eq:hilbert-skew-adjoint}
\end{equation}
Hilbert 变换因此是反自伴算符，而 $\ii\mathcal H$ 是自伴算符。

由乘子再作用一次，
\[
\widehat{\mathcal H^2 f}(k)
=(-\ii)^2\operatorname{sgn}^2(k)\widehat f(k)
=-\widehat f(k)
\]
对几乎所有 $k$ 成立。单点 $k=0$ 对 $L^2$ 函数没有影响，所以在整个 $L^2(\mathbb R)$ 上
\begin{equation}
\boxed{\mathcal H^2=-I.}
\label{eq:hilbert-square}
\end{equation}
结合 \eqnrefs{eq:hilbert-skew-adjoint,eq:hilbert-square}，还可核对
\[
\mathcal H^\dagger\mathcal H=(-\mathcal H)\mathcal H=I,
\]
即 $\mathcal H$ 同时是一个酉算符。

这一结构立刻给出正、负频率投影。定义
\begin{equation*}P_+=\frac12(I+\ii\mathcal H),
\qquad
P_-=\frac12(I-\ii\mathcal H).\end{equation*}
在 Fourier 空间，$\ii\mathcal H$ 的乘子就是 $\operatorname{sgn}(k)$，因此
\[
\widehat{P_+f}(k)
=\frac{1+\operatorname{sgn}(k)}2\widehat f(k),
\qquad
\widehat{P_-f}(k)
=\frac{1-\operatorname{sgn}(k)}2\widehat f(k).
\]
除去测度为零的 $k=0$，前者保留 $k>0$ 的频率，后者保留 $k<0$ 的频率。由 $\mathcal H^2=-I$ 可直接算得
\[
P_+^2=P_+,
\qquad
P_-^2=P_-,
\qquad
P_+P_-=0,
\qquad
P_++P_-=I.
\]
所以 Hilbert 变换不是孤立出现的奇异积分技巧，它实际上编码了 $L^2(\mathbb R)$ 中的频率分解。

最后说明几个以后经常使用、但容易凭感觉写错的交换关系。由 Fourier 乘子可见，只要各项属于相应函数空间，
\[
\mathcal H(\partial_x f)=\partial_x(\mathcal H f),
\]
因为两边的 Fourier 乘子都是
\[
(-\ii\operatorname{sgn}k)(\ii k)=|k|.
\]
同样地，平移算符 $\tau_a f(x)=f(x-a)$ 与 $\mathcal H$ 对易，因为平移只带来相位 $\ee^{-\ii ka}$；对正尺度 $a>0$，若 $D_af(x)=f(ax)$，也有
\[
\mathcal H D_a=D_a\mathcal H.
\]
这些关系解释了为什么 Hilbert 变换在解析信号、色散方程和奇异积分方程中反复出现：它既保持 $L^2$ 能量，又与最基本的平移和微分结构相容。

课堂记录还给出了形式

\[
y(x)=\int_C(x-t)^\mu v(t)\,\dd t.
\]

仅凭这个式子还不能唯一确定一个具体变换：必须同时指定积分路径 $C$、复幂 $(x-t)^\mu$ 的分支、端点行为以及使积分收敛的参数范围。若 $C$ 取实区间并固定适当分支，这类核与 Abel 积分以及 Riemann--Liouville 型分数积分有直接联系；若 $C$ 是复平面路径，则分支切割和绕行方式会改变结果。因此在这些条件给出之前，只把它保留为一类形式积分变换，而不强行写出不存在唯一依据的逆公式。
